% Diagramma TikZ a due corsie per l'asimmetria di campionamento e la sparsita' strutturale
% Allineato graficamente e concettualmente alla Slide 03/04 della presentazione
\begin{tikzpicture}[
    font=\sffamily,
    >=Stealth,
    axisLine/.style={->, line width=0.9pt, color=brandDark},
    gridLine/.style={color=brandBorderStrong, line width=0.5pt, dashed},
    traceHR/.style={color=brandPrimary, line width=1.3pt, line join=round, line cap=round},
    traceBP/.style={color=brandCoral, line width=1.4pt, line cap=round},
    voidLine/.style={color=brandCoral!60, line width=1.1pt, dashed},
    spanArrow/.style={<->, line width=1.1pt, color=brandCoral}
]

    % ========================================================
    % CORSIA 1: FREQUENZA CARDIACA (BPM / ECG) - FLUSSO CONTINUO A 1 Hz
    % ========================================================
    \draw[fill=brandSurface, draw=brandBorderStrong, rounded corners=6pt, line width=0.8pt] 
        (-0.2, 3.5) rectangle (12.8, 5.7);
    
    % Intestazione corsia 1
    \node[anchor=west, font=\scriptsize\bfseries\color{brandDark}] at (0.0, 5.42) 
        {FREQUENZA CARDIACA (HR / BPM) $\cdot$ FLUSSO CONTINUO};


    % Riquadro bianco interno segnale (stile monitor/presentazione)
    \draw[fill=white, draw=brandBorder, rounded corners=4pt] (0.0, 3.7) rectangle (12.6, 5.1);

    % Griglia minuti sincronizzata (ogni 2.4 cm = 1 minuto / 60s, totale 300s = 5 min = 12.0 cm)
    \foreach \t/\x in {0/0.3, 60/2.7, 120/5.1, 180/7.5, 240/9.9, 300/12.3} {
        \draw[gridLine] (\x, 3.7) -- (\x, 5.1);
    }

    % Linea isoelettrica di riferimento
    \draw[color=brandPrimary!12, line width=0.6pt] (0.0, 4.35) -- (12.6, 4.35);

    % Onda classica continua BPM/ECG (complessi QRS regolari e onde P-T)
    % Ciascun ciclo cardiaco ha passo 0.6cm (21 battiti rappresentati da x=0.3 a x=12.3)
    \draw[traceHR] (0.0, 4.35) -- (0.3, 4.35)
        % Ciclo 1
        -- ++(0.12, 0) to[out=60,in=120] ++(0.08, 0) -- ++(0.06, 0) -- ++(0.03, -0.06) -- ++(0.05, 0.52) -- ++(0.05, -0.66) -- ++(0.03, 0.20) -- ++(0.05, 0) to[out=50,in=130] ++(0.10, 0) -- ++(0.03, 0)
        % Ciclo 2
        -- ++(0.12, 0) to[out=60,in=120] ++(0.08, 0) -- ++(0.06, 0) -- ++(0.03, -0.06) -- ++(0.05, 0.52) -- ++(0.05, -0.66) -- ++(0.03, 0.20) -- ++(0.05, 0) to[out=50,in=130] ++(0.10, 0) -- ++(0.03, 0)
        % Ciclo 3
        -- ++(0.12, 0) to[out=60,in=120] ++(0.08, 0) -- ++(0.06, 0) -- ++(0.03, -0.06) -- ++(0.05, 0.52) -- ++(0.05, -0.66) -- ++(0.03, 0.20) -- ++(0.05, 0) to[out=50,in=130] ++(0.10, 0) -- ++(0.03, 0)
        % Ciclo 4
        -- ++(0.12, 0) to[out=60,in=120] ++(0.08, 0) -- ++(0.06, 0) -- ++(0.03, -0.06) -- ++(0.05, 0.52) -- ++(0.05, -0.66) -- ++(0.03, 0.20) -- ++(0.05, 0) to[out=50,in=130] ++(0.10, 0) -- ++(0.03, 0)
        % Ciclo 5
        -- ++(0.12, 0) to[out=60,in=120] ++(0.08, 0) -- ++(0.06, 0) -- ++(0.03, -0.06) -- ++(0.05, 0.52) -- ++(0.05, -0.66) -- ++(0.03, 0.20) -- ++(0.05, 0) to[out=50,in=130] ++(0.10, 0) -- ++(0.03, 0)
        % Ciclo 6
        -- ++(0.12, 0) to[out=60,in=120] ++(0.08, 0) -- ++(0.06, 0) -- ++(0.03, -0.06) -- ++(0.05, 0.52) -- ++(0.05, -0.66) -- ++(0.03, 0.20) -- ++(0.05, 0) to[out=50,in=130] ++(0.10, 0) -- ++(0.03, 0)
        % Ciclo 7
        -- ++(0.12, 0) to[out=60,in=120] ++(0.08, 0) -- ++(0.06, 0) -- ++(0.03, -0.06) -- ++(0.05, 0.52) -- ++(0.05, -0.66) -- ++(0.03, 0.20) -- ++(0.05, 0) to[out=50,in=130] ++(0.10, 0) -- ++(0.03, 0)
        % Ciclo 8
        -- ++(0.12, 0) to[out=60,in=120] ++(0.08, 0) -- ++(0.06, 0) -- ++(0.03, -0.06) -- ++(0.05, 0.52) -- ++(0.05, -0.66) -- ++(0.03, 0.20) -- ++(0.05, 0) to[out=50,in=130] ++(0.10, 0) -- ++(0.03, 0)
        % Ciclo 9
        -- ++(0.12, 0) to[out=60,in=120] ++(0.08, 0) -- ++(0.06, 0) -- ++(0.03, -0.06) -- ++(0.05, 0.52) -- ++(0.05, -0.66) -- ++(0.03, 0.20) -- ++(0.05, 0) to[out=50,in=130] ++(0.10, 0) -- ++(0.03, 0)
        % Ciclo 10
        -- ++(0.12, 0) to[out=60,in=120] ++(0.08, 0) -- ++(0.06, 0) -- ++(0.03, -0.06) -- ++(0.05, 0.52) -- ++(0.05, -0.66) -- ++(0.03, 0.20) -- ++(0.05, 0) to[out=50,in=130] ++(0.10, 0) -- ++(0.03, 0)
        % Ciclo 11
        -- ++(0.12, 0) to[out=60,in=120] ++(0.08, 0) -- ++(0.06, 0) -- ++(0.03, -0.06) -- ++(0.05, 0.52) -- ++(0.05, -0.66) -- ++(0.03, 0.20) -- ++(0.05, 0) to[out=50,in=130] ++(0.10, 0) -- ++(0.03, 0)
        % Ciclo 12
        -- ++(0.12, 0) to[out=60,in=120] ++(0.08, 0) -- ++(0.06, 0) -- ++(0.03, -0.06) -- ++(0.05, 0.52) -- ++(0.05, -0.66) -- ++(0.03, 0.20) -- ++(0.05, 0) to[out=50,in=130] ++(0.10, 0) -- ++(0.03, 0)
        % Ciclo 13
        -- ++(0.12, 0) to[out=60,in=120] ++(0.08, 0) -- ++(0.06, 0) -- ++(0.03, -0.06) -- ++(0.05, 0.52) -- ++(0.05, -0.66) -- ++(0.03, 0.20) -- ++(0.05, 0) to[out=50,in=130] ++(0.10, 0) -- ++(0.03, 0)
        % Ciclo 14
        -- ++(0.12, 0) to[out=60,in=120] ++(0.08, 0) -- ++(0.06, 0) -- ++(0.03, -0.06) -- ++(0.05, 0.52) -- ++(0.05, -0.66) -- ++(0.03, 0.20) -- ++(0.05, 0) to[out=50,in=130] ++(0.10, 0) -- ++(0.03, 0)
        % Ciclo 15
        -- ++(0.12, 0) to[out=60,in=120] ++(0.08, 0) -- ++(0.06, 0) -- ++(0.03, -0.06) -- ++(0.05, 0.52) -- ++(0.05, -0.66) -- ++(0.03, 0.20) -- ++(0.05, 0) to[out=50,in=130] ++(0.10, 0) -- ++(0.03, 0)
        % Ciclo 16
        -- ++(0.12, 0) to[out=60,in=120] ++(0.08, 0) -- ++(0.06, 0) -- ++(0.03, -0.06) -- ++(0.05, 0.52) -- ++(0.05, -0.66) -- ++(0.03, 0.20) -- ++(0.05, 0) to[out=50,in=130] ++(0.10, 0) -- ++(0.03, 0)
        % Ciclo 17
        -- ++(0.12, 0) to[out=60,in=120] ++(0.08, 0) -- ++(0.06, 0) -- ++(0.03, -0.06) -- ++(0.05, 0.52) -- ++(0.05, -0.66) -- ++(0.03, 0.20) -- ++(0.05, 0) to[out=50,in=130] ++(0.10, 0) -- ++(0.03, 0)
        % Ciclo 18
        -- ++(0.12, 0) to[out=60,in=120] ++(0.08, 0) -- ++(0.06, 0) -- ++(0.03, -0.06) -- ++(0.05, 0.52) -- ++(0.05, -0.66) -- ++(0.03, 0.20) -- ++(0.05, 0) to[out=50,in=130] ++(0.10, 0) -- ++(0.03, 0)
        % Ciclo 19
        -- ++(0.12, 0) to[out=60,in=120] ++(0.08, 0) -- ++(0.06, 0) -- ++(0.03, -0.06) -- ++(0.05, 0.52) -- ++(0.05, -0.66) -- ++(0.03, 0.20) -- ++(0.05, 0) to[out=50,in=130] ++(0.10, 0) -- ++(0.03, 0)
        % Ciclo 20
        -- ++(0.12, 0) to[out=60,in=120] ++(0.08, 0) -- ++(0.06, 0) -- ++(0.03, -0.06) -- ++(0.05, 0.52) -- ++(0.05, -0.66) -- ++(0.03, 0.20) -- ++(0.05, 0) to[out=50,in=130] ++(0.10, 0) -- ++(0.03, 0)
        -- (12.6, 4.35);


    % ========================================================
    % CORSIA 2: PRESSIONE ARTERIOSA (NIBP) - MISURE DISCONTINUE EQUIDISTANTI
    % ========================================================
    \draw[fill=brandSurface, draw=brandBorderStrong, rounded corners=6pt, line width=0.8pt] 
        (-0.2, 0.8) rectangle (12.8, 3.2);

    % Intestazione corsia 2
    \node[anchor=west, font=\scriptsize\bfseries\color{brandDark}] at (0.0, 2.92) 
        {PRESSIONE ARTERIOSA (BRACCIALE NIBP) $\cdot$ MISURAZIONI DISCONTINUE};


    % Riquadro bianco interno segnale
    \draw[fill=white, draw=brandBorder, rounded corners=4pt] (0.0, 1.0) rectangle (12.6, 2.6);

    % Griglia minuti sincronizzata
    \foreach \t/\x in {0/0.3, 60/2.7, 120/5.1, 180/7.5, 240/9.9, 300/12.3} {
        \draw[gridLine] (\x, 1.0) -- (\x, 2.6);
    }

    % --- 1° MISURAZIONE A t = 0 min (x = 0.3): treno di 2 battiti pressori oscillometrici ---
    \draw[traceBP] (0.30, 1.30) 
        .. controls (0.35, 1.30) and (0.40, 2.38) .. (0.45, 2.45)
        .. controls (0.49, 2.50) and (0.54, 1.82) .. (0.58, 1.68)
        .. controls (0.60, 1.62) and (0.64, 1.80) .. (0.68, 1.65)
        .. controls (0.74, 1.48) and (0.80, 1.30) .. (0.85, 1.30)
        % Battito 2
        .. controls (0.90, 1.30) and (0.95, 2.38) .. (1.00, 2.45)
        .. controls (1.04, 2.50) and (1.09, 1.82) .. (1.13, 1.68)
        .. controls (1.15, 1.62) and (1.19, 1.80) .. (1.23, 1.65)
        .. controls (1.29, 1.48) and (1.35, 1.30) .. (1.40, 1.30);
    \node[above, font=\tiny\bfseries\color{brandCoral}] at (0.85, 2.45) {118/76\,mmHg};
    \node[below, font=\tiny\bfseries\color{brandDark}, fill=brandSurface, rounded corners=2pt, inner sep=1pt] at (0.85, 1.25) {1° Gonfiaggio};

    % --- ZONA DI VUOTO STRUTTURALE 1 (4 MINUTI EQUIDISTANTI: x = 1.40 -> x = 9.40) ---
    \draw[voidLine] (1.40, 1.30) -- (9.40, 1.30);



    % --- 2° MISURAZIONE EQUIDISTANTE A t = 4 min (x = 9.9 - 0.5 = 9.4): treno di battiti pressori ---
    \draw[traceBP] (9.40, 1.30) 
        .. controls (9.45, 1.30) and (9.50, 2.38) .. (9.55, 2.45)
        .. controls (9.59, 2.50) and (9.64, 1.82) .. (9.68, 1.68)
        .. controls (9.70, 1.62) and (9.74, 1.80) .. (9.78, 1.65)
        .. controls (9.84, 1.48) and (9.90, 1.30) .. (9.95, 1.30)
        % Battito 2
        .. controls (10.00, 1.30) and (10.05, 2.38) .. (10.10, 2.45)
        .. controls (10.14, 2.50) and (10.19, 1.82) .. (10.23, 1.68)
        .. controls (10.25, 1.62) and (10.29, 1.80) .. (10.33, 1.65)
        .. controls (10.39, 1.48) and (10.45, 1.30) .. (10.50, 1.30);
    \node[above, font=\tiny\bfseries\color{brandCoral}] at (9.95, 2.45) {122/79\,mmHg};
    \node[below, font=\tiny\bfseries\color{brandDark}, fill=brandSurface, rounded corners=2pt, inner sep=1pt] at (9.95, 1.25) {2° Gonfiaggio};

    % Coda vuota verso fine finestra (t = 5 min)
    \draw[voidLine] (10.50, 1.30) -- (12.6, 1.30);


    % ========================================================
    % ASSE TEMPORALE COMUNE EQUIDISTANTE E SINCRONIZZATO (0 - 5 MINUTI)
    % ========================================================
    \draw[axisLine] (-0.2, 0.4) -- (12.8, 0.4) node[right, font=\scriptsize\bfseries] {$t$};
    
    \foreach \min/\sec/\x in {0\,min/0\,s/0.3, 1\,min/60\,s/2.7, 2\,min/120\,s/5.1, 3\,min/180\,s/7.5, 4\,min/240\,s/9.9, 5\,min/300\,s/12.3} {
        \draw[line width=0.9pt, color=brandDark] (\x, 0.4) -- (\x, 0.26) 
            node[below=2pt, font=\scriptsize\bfseries\color{brandDark}] {\min}
            node[below=11pt, font=\tiny\color{brandMuted}] {(\sec)};
    }



\end{tikzpicture}
